\documentclass{article}
\usepackage{amsmath}
\usepackage{amssymb}

\usepackage{graphicx}

% Encodage et langue (optionnel mais recommandé)
\usepackage[utf8]{inputenc}
\usepackage[T1]{fontenc}
\usepackage[french]{babel}

\begin{document}

\title{Lot D.1 : résolution de probleme d'optimisation}
\author{Daouda CISSE : innovation,
Hicham El KHAZROUNI : Experimentation ,
Nael CHANEGRIH : résultat}
\date{}
\maketitle

\section*{Problème à résoudre}
Maximiser la distance minimale entre une particule et les autres dans un espace 2D
de taille $h \times w$, après $t$ secondes, en fixant leur vitesse et leur position initiales.
On note \textit{optimum} cette valeur optimale dans la suite.

\section*{Description de l'innovation}
Nous faisons les hypothèses suivantes, qui seront justifiées un peu plus loin :
$n$ est petit, de l’ordre de 5.

On choisit de donner aux particules une vitesse initiale de sorte
qu’elles soient, à tout instant, uniformément réparties sur un cercle, et qu’à
$t$ secondes elles se trouvent sur le cercle de rayon $\min(h,w)$.

\section*{Explication des choix, des hypothèses et des inconvénients}
On fixe une borne :

Dans un espace de surface $h^2$ avec $n$ particules,
la densité surfacique par particule est $\frac{h^2}{n}$.
On en déduit que la distance minimale maximale
entre deux particules ne peut dépasser
$\sqrt{\frac{h^2}{n}} = \frac{h}{\sqrt{n}}$.

\subsection*{Alternative non choisie}
Si $n$ est environ égal à $m^2$, on peut construire une grille uniforme de taille $m \times m$
que l’on agrandit autant que possible. Alors, si $h = w$, on a :
\[
(k - 1)\,\textit{optimum} = h
\]
d’où une valeur de l’optimum très proche de $\frac{h}{\sqrt{n}}$, c’est-à-dire de la borne.

Cependant, il est très compliqué d’initialiser les particules (en raison des collisions et d’autres facteurs)
pour obtenir ce résultat. C’est pourquoi nous avons choisi une répartition circulaire plutôt qu’une grille.

\subsection*{Alternative choisie}
Notre méthode aboutit à des particules réparties sur un cercle de rayon
\[
\text{ray} = \frac{\min(h,w)}{2}.
\]
La distance minimale maximale entre deux particules est alors
\[
\frac{2\pi \cdot \text{ray}}{n} = \frac{\pi \cdot \min(h,w)}{n}.
\]

Pour $h = w$, on obtient une valeur approximativement égale à $\frac{3h}{n}$,
ce qui est proche de l’optimum pour des valeurs faibles de $n$.

Cependant, pour des valeurs très grandes de $n$, on s’en éloigne de manière quadratique,
car le dénominateur est en $n$ et non en $\sqrt{n}$, d’où l’hypothèse que $n$ est faible
